%% Auto-generated by code/make_text.py from code/results/*.json. Do not edit by hand.
\newcommand{\RESabstract}{over all 15 model--input combinations, test accuracy ranges from 9.8--22.8\% and macro-F1 from 2.7--14.4\%, against 11.1\% accuracy for always predicting the most frequent class and 10.0\% for random guessing; the best configuration selected on the dev split (CLIP ViT-L/14, dual-pol input) reaches 22.8\% accuracy and 14.4\% macro-F1}
\newcommand{\RESprobe}{59.3--70.1\% accuracy from only 5,400 labeled images}
\newcommand{\RESpromptsText}{Accuracy ranges from 9.0\% (\texttt{landcover}) to 12.0\% (\texttt{structure}), a spread of 3.0 percentage points around the majority-class baseline of 11.1\%. Macro-F1 lies between 2.3--6.3\%, below the $\approx$10\% expected from random guessing, and balanced accuracy between 8.3--10.8\%. With 21,600 test images the 95\% confidence intervals are about $\pm$0.4 points wide, so some differences between templates are statistically detectable, but none is practically meaningful: no template brings CLIP above chance level on SAR. The spatial template (``an aerial view of'', 11.6\%) and the three SAR-specific templates (9.8--11.4\%) perform alike, so the data do not support a claim that spatial prompts substantially outperform radar terminology. Several templates never predict some classes at all (8--10 of 10 classes are ever predicted).}
\newcommand{\RESmodelsText}{Table~\ref{tab:models} and Fig.~\ref{fig:modes} compare all models and inputs. Every configuration stays close to chance: accuracy 9.8--22.8\%, macro-F1 2.7--14.4\%. With VV input the highest macro-F1 is obtained by GeoRSCLIP ViT-B/32 (12.8\%). Optical remote-sensing pre-training does not reliably help or hurt: RemoteCLIP ViT-B/32 versus CLIP ViT-B/32: 12.0\% vs.\ 9.8\% accuracy and 6.1\% vs.\ 6.3\% macro-F1 (McNemar $p<0.001$); RemoteCLIP ViT-L/14 versus CLIP ViT-L/14: 12.0\% vs.\ 12.7\% accuracy and 4.7\% vs.\ 10.3\% macro-F1 (McNemar $p=0.022$). Domain-specific pre-training on optical imagery therefore does not translate into a consistent advantage on SAR; the differences are small in either direction. The larger ViT-L/14 backbones do not change this picture (CLIP ViT-L/14: 10.3, RemoteCLIP ViT-L/14: 4.7\% macro-F1 with VV). The effect of polarization depends on the model: the dual-pol composite raises macro-F1 for CLIP ViT-B/32, CLIP ViT-L/14 (+5.0, +4.1 points) and lowers it for RemoteCLIP ViT-B/32, GeoRSCLIP ViT-B/32 (-3.4, -9.8 points). A false-color composite therefore gives the image encoder more information but also moves the input further from the images these models were trained on.}
\newcommand{\REStechText}{Table~\ref{tab:techniques} shows that none of the standard zero-shot techniques closes the gap. Relative to the single dev-selected template, accuracy changes by prompt ensembling -5.4 to +0.9; physics-informed descriptions -1.1 to +4.5; rotation TTA -1.8 to +0.4; the combination of all three -2.5 to -0.4 percentage points across models, and all results remain between 8.0--21.2\%. The physics-informed descriptions, which encode expected radar behavior (e.g.\ volume scattering in forests, specular reflection on water), do not help the text encoders, which were never trained to associate such descriptions with radar images.}
\newcommand{\REScoarse}{coarse accuracy exceeds the majority-class baseline by at most 1.3 points for the 3-group hierarchy and exceeds the majority-class baseline by at most 6.8 points for the 5-group hierarchy}
\newcommand{\REShierText}{Coarse accuracies must be read against the majority-class baseline, which is much higher than at the fine level because the groups are unequal in size. For H3, always predicting ``vegetation'' already gives 50.0\% accuracy; the best zero-shot coarse result is 51.3\% (RemoteCLIP ViT-L/14, fine$\to$coarse), and 1 of 10 configurations exceed the baseline. Balanced accuracy (28.5--44.4\%) stays close to the 33.3\% of random guessing. For H5 the majority baseline is 29.6\%, the best result 36.4\%, and 7 of 10 configurations exceed it. The large jump from fine to coarse accuracy is therefore mostly due to having fewer, larger classes rather than to better recognition. A classifier that assigns every image to the largest group reaches exactly this baseline, so coarse accuracy alone, without the baseline, can be badly misleading.}
\newcommand{\RESperclassText}{Per-class results (Table~\ref{tab:perclass}, Fig.~\ref{fig:confusion}) show how the errors arise. For CLIP ViT-B/32 with VV input, 32.9\% of all test images are assigned to ``Forest'', and recall is low for every class (at most 38.5\%). In the best zero-shot configuration the highest F1 scores are obtained for SeaLake (74.7\%), Forest (28.7\%) and AnnualCrop (24.2\%), while 6 classes (HerbaceousVegetation, Highway, Industrial, Pasture, Residential, River) have F1 below 5\%. Large, calm water bodies (SeaLake) fill the whole patch with very low, homogeneous backscatter (Fig.~\ref{fig:examples}) and are the most distinctive class; narrow rivers occupy only part of a patch and are not recognized. Most other classes, including the urban classes with bright double-bounce returns, are recognized poorly, which suggests that the text embeddings do not capture these SAR signatures.}
\newcommand{\RESgapText}{Table~\ref{tab:gap} locates the failure. As a positive control, the identical pipeline applied to optical EuroSAT reaches 30.7--45.7\% zero-shot accuracy with RGB input and 20.8--33.3\% with grayscale input (template ``an aerial view of''), confirming that the implementation works and that removing color costs part of the performance. On SAR the same template gives 8.3--11.6\%. A logistic-regression linear probe trained on the 5,400 dev images (not zero-shot) reaches 59.3--65.8\% on VV and 63.7--70.1\% on dual-pol input, far above chance. The frozen image encoders therefore do produce SAR embeddings that separate the land-cover classes; what fails is the match between these embeddings and the text embeddings of the class names. SAR images look unlike anything the text encoders were aligned with during pre-training, so prompt wording alone cannot recover the class information that the image embeddings contain.}
\newcommand{\RESpitfallText}{SAR products such as EuroSAT-SAR store backscatter in decibels, and almost every value is negative. If the floating-point array is cast directly to 8-bit integers (for example by clipping to 0--255 before saving or displaying it), every pixel becomes 0 and all images are black. We reproduced this error: the ViT-B/32 models then predict only 2--4 distinct classes, with accuracies of 7.6--12.1\% and macro-F1 of 1.6--3.4\%. Such results look stable and are close to the majority-class rate, so they can easily be mistaken for a weak but genuine baseline. Checking the number of distinct predicted classes and comparing against a majority-class baseline exposes the problem immediately; we recommend both as routine checks for SAR experiments with VLMs.}
\newcommand{\RESdiscussionText}{Our results do not support the use of current CLIP-style VLMs as zero-shot classifiers for single-date Sentinel-1 land cover. Fine-grained zero-shot accuracy is close to chance for every model, input and prompt we tested; coarse accuracy exceeds the majority-class baseline by at most 1.3 points for the 3-group hierarchy and exceeds the majority-class baseline by at most 6.8 points for the 5-group hierarchy; and the best configuration (22.8\% accuracy) would not be useful in practice. The linear-probe results suggest that the more promising direction is to keep the frozen image encoder and learn the class representations from a few labeled SAR examples, for instance through prompt learning~\cite{zhou2022learning} or SAR--text pre-training. \textit{Limitations.} (i) We evaluate one dataset (EuroSAT-SAR): European sites, $64\times64$ pixel patches at 10\,m, a single Sentinel-1 acquisition mode and season per patch, and 10 classes inherited from optical labels; results may differ for other regions, higher-resolution SAR, time series or other class definitions. (ii) The VLMs require three-channel input. Duplicating VV or VH, or building a false-color composite, are pragmatic mappings that do not match the statistics of the images the models were trained on; this modality mismatch is one of the causes of the poor zero-shot performance. (iii) We test five dual-encoder models and hand-written prompts; SAR-specific foundation models, generative VLMs and learned prompts are outside the scope of this study. (iv) Rotation TTA and the optical positive control were run for the ViT-B/32 models only.}
\newcommand{\RESconclusionText}{This study provides a reproducible baseline for zero-shot vision-language classification of real dual-polarization Sentinel-1 imagery. Across five frozen VLMs, ten prompt templates, three input encodings and four standard zero-shot techniques, fine-grained accuracy on EuroSAT-SAR stays close to chance (9.8--22.8\%, majority baseline 11.1\%). Coarse hierarchical classification is no remedy: coarse accuracy exceeds the majority-class baseline by at most 1.3 points for the 3-group hierarchy and exceeds the majority-class baseline by at most 6.8 points for the 5-group hierarchy. Prompt wording matters little, and optical remote-sensing pre-training does not transfer reliably to SAR. The image embeddings nevertheless contain land-cover information (linear probe 59.3--70.1\%), which identifies image--text alignment as the bottleneck. We therefore caution against using general-purpose or optical remote-sensing VLMs as zero-shot SAR classifiers, and we see few-shot adaptation and SAR-specific vision-language pre-training as the necessary next steps. Finally, we recommend reporting majority-class baselines, class-balanced metrics and the number of predicted classes, which together reveal degenerate results such as those caused by decibel-to-8-bit conversion.}
